\subsection{Primary terminal intervals with their regular entrances}
\label{rw:anchors:primary}

Let $R$ be the regular prefix of the physical interval and $T$ its
terminal union. On their actual open overlap, regular changes of input
identify the same original-central $K$.

By Proposition~\ref{rw:regular-chain}, $R$ has at most one simple
zero. After a zero, or under the signed entrance below, it supplies
the departure sign $-\sigma_1$ and the coefficient cone
\eqref{rw:merge:primary-cone}. More generally,
Lemma~\ref{rw:finite:regular-direction} gives that derivative sign
whenever $K$ has sign $-\sigma_1$, without requiring the cone;
if $\sigma_1=0$, $K,K_p$ have sign $-\tau_1$ off the center.
A regular zero requires $\tau_1\sigma_1<0$.
These statements hold on the actual overlap, including nonclosing inputs.

\label{rw:anchors:seam-convention}
If $K$ is constant on a terminal interval, a nonzero constant
contributes no zero; a zero constant extends into its regular overlap
and, by the regular count under forcing, belongs to the exact center
case, contributing no isolated closing zero. Otherwise choose a seam
in the overlap away from zeros of $K$ and the terminal derivative
numerator, after a regular zero if one has occurred. The part of $T$ before the seam is then assigned to $R$,
with every multiplicity retained. We fix this seam before classifying
the terminal critical zeros, and use the same convention for the
secondary anchors and the final union.

\begin{lemma}[Entrance counts]\label{rw:anchors:calculus}
On a tail after this seam, a positive-multiplier numerator with at
most $j$ zeros gives $Z(K)\le j+1$, or $Z(K)\le j$ when
$K K'_{\rm out}>0$ at the seam. On an explicitly constructed
comparison continuation, a two-zero bound and opposite nonzero
endpoint signs give at most one zero on every physical subinterval.
If $u$ decreases outward and a positive-multiplier numerator of
$K_u$ strictly increases in $u$, an entrance with $K<0$ and
$K'_{\rm out}<0$ likewise leaves at most one zero, which is simple.
All counts include multiplicity and apply to finite selections at
excluded endpoints.
\end{lemma}
\begin{proof}
Ordinary and away Rolle in Lemma~\ref{rw:counting:rolle} give the
first assertion. Opposite endpoint signs force odd total multiplicity,
so a two-zero continuation has one zero. In the last assertion the
seam numerator is positive and, moving outward, can become negative
at most once. Thus $K$ first decreases below zero and then strictly
increases, allowing at most one simple crossing.
\end{proof}

\begin{proposition}[Regular lower arcs at an upper saddle]
\label{rw:anchors:regular-saddle}
In the chart of Proposition~\ref{rw:finite:regular-saddle}, the terminal
part has at most two zeros of $K$, with multiplicity. If $R$ has a zero,
its remaining terminal part has at most one. The same improvement holds
when $R$ has no zero but the supplied entrance satisfies
\begin{equation}\label{rw:anchors:signed-entry}
 \tau_1\sigma_1<0,\qquad
 \operatorname{sign}K=-\operatorname{sign}\sigma_1.
\end{equation}
Thus $Z(K;R\cup T)\le2$ without a prescribed entrance and it is
at most one under \eqref{rw:anchors:signed-entry} if the preceding
part has no zero.
\end{proposition}
\begin{proof}
Proposition~\ref{rw:finite:regular-saddle} gives the unanchored bound.
In either improved case, the regular overlap supplies
\eqref{rw:anchors:signed-entry} and the cone
\eqref{rw:merge:primary-cone}; coefficient transport supplies
\eqref{rw:finite:saddle-relative-forcing} on this fixed fibre.
Lemma~\ref{rw:finite:saddle-entry-signs} then gives opposite signs
at the seam and $t=0$ for the same original-central $K$ on its
two-zero comparison continuation. Lemma~\ref{rw:anchors:calculus}
leaves one zero there and hence on its physical part. Add the
possible single regular zero.
\end{proof}

At the three-saddle end the derivative numerator may have two zeros.
Ordinary Rolle would then allow three zeros of $K$. We need the signs
forced by the two-critical-zero case, together with the regular entrance,
to remove this extra possibility.

\begin{proposition}[The three-saddle union at a fixed seam]
\label{rw:anchors:three-saddle}
Suppose $T$ is the three-saddle chart of
Proposition~\ref{rw:finite:three-count}, and impose its relative
forcing on the fixed fibre. Then
\[
 Z(K;R\cup T)\le2.
\]
After a zero in $R$, every seam satisfying the preceding conventions
leaves at most one terminal zero. If there is no regular zero and
\eqref{rw:anchors:signed-entry} holds, the union has at most one zero.
All counts include multiplicity and remain valid when central positivity
or a moving lower boundary truncates the passing interval.
\end{proposition}
\begin{proof}
For the physical interval $J=[p_0,p_\gamma]$ of
Proposition~\ref{rw:saturation}, the outer closing orbit supplies the
local forcing in \eqref{rw:finite:three-forcing}.
Indeed, write the actual seam as $p_s=P_s-\psi_+(t_s)$,
with $\psi_+'>0$. The separatrix input $P_s$ has no finite upper
return. Saturation therefore excludes $p_\gamma\ge P_s$, and hence
\[
 p_s\le p_\gamma<P_s,\qquad 0<t_\gamma\le t_s.
\]
The fixed signed connecting maps constructed in
Section~\ref{rw:finite:three-saddle} contain this smaller upper
passage. Its actual lower returns remain on the passing sides
of the two strong separatrices, by
Proposition~\ref{rw:lower-preparation}; thus
$G_+(t_\gamma)>0$ and $G_-(v(t_\gamma))>0$.
Consequently $\gamma$ is an actual closing zero in the same
three-saddle comparison chart. Equation~\eqref{rw:finite:three-forcing}
holds on its fixed coefficient fibre, including the comparison
continuation used below.

Write $u=G_+(t)$, $z=G_-(v(t))$, $v=t^re^w$ as in
Section~\ref{rw:finite:three-saddle}. On the deficit comparison domain,
\[
 M=tuz K_t,\qquad t>0,\quad u>0,\quad z>0.
\]
The outward coordinate is $p=P_s-\psi_+(t)$ with
$\psi_+'>0$. Let $T_s$ be the actual positive-central tail after
the fixed seam $t_s$, and set
\begin{equation}\label{rw:anchors:fixed-seam-number}
 j_s=Z(M;T_s).
\end{equation}
The local bound gives $j_s\le2$ at this fixed seam.
If $j_s\le1$, Lemma~\ref{rw:anchors:calculus} proves the assertions:
regular departure gives $KK_p>0$ after a regular zero or under the
prescribed signed entry.

Suppose now that $j_s=2$. Apply, if needed, the reflection identity to
the whole union, its coefficients, and its already fixed seam so that
$r>1$. It sends $K$ to its negative composed with the corresponding
input change, and changes the signs of $\tau_1,\sigma_1$ together.
It preserves multiplicity, outward ordering and the signed-entry
condition, and transports the already chosen seam and threshold to
their corresponding actual points. The reflected outgoing label is
the original incoming label; the numerical outgoing cutoff need not
be unchanged.

With this local forcing fixed, apply
Proposition~\ref{rw:finite:three-count} on the same deficit comparison
interval, with the original central function $K=\mathcal K$ and
the already fixed seam. Its two-zero classification gives
\begin{equation}\label{rw:anchors:three-end-signs}
 u_0,z_0>0,\qquad \tau_1>0,\qquad
 M(0+)=\epsilon\kappa u_0z_0>0,\qquad K(0+)=-\infty.
\end{equation}
The actual tail's two zeros exhaust the comparison bound.
The same proposition gives $M(t_s)>0$ and $K_p(t_s)<0$ at its
fixed seam.

The regular overlap determines the remaining sign, $K(t_s)$.
For $\sigma_1\ge0$, $\tau_1>0$ and the positive regular columns
give $K(t_s)<0$, with relative forcing absorbed. For $\sigma_1<0$,
a positive value of $K(t_s)$ would have sign $-\sigma_1$, and
Lemma~\ref{rw:finite:regular-direction} would imply $K_p(t_s)>0$,
a contradiction. The seam avoids a zero, so again $K(t_s)<0$.
Together with \eqref{rw:anchors:three-end-signs}, this gives negative
signs at both ends of the continued tail. Ordinary Rolle bounds its
total multiplicity by three, and parity reduces this to two.
Restricting to the physical subinterval preserves that bound.

There cannot have been a regular zero: since $\tau_1>0$, such a zero
would require $\sigma_1<0$ and would leave $K>0$ at the seam,
contrary to the preceding conclusion. Thus the union has at most two
zeros. Finally the no-zero signed-entry alternative also excludes
$j_s=2$. If $\sigma_1>0$, it contradicts $\tau_1\sigma_1<0$;
if $\sigma_1<0$, it would require a positive seam value of $K$.
Only $j_s\le1$ remains, and the away estimate gives one.
\end{proof}

\begin{proposition}[Primary weak thresholds with their regular part]
\label{rw:anchors:primary-weak}
Let $T$ be either a positive-$\bar B$ weak threshold with a lower
root at $P$, as in Proposition~\ref{rw:infinity:weak-loss}, or the weak-center chart of
Proposition~\ref{rw:infinity:weak-center}. Under the actual overlap
and fixed-fibre forcing hypotheses above,
\[
 Z(K;R\cup T)\le2.
\]
After a zero in $R$, the remaining terminal part has at most one zero.
With no regular zero and \eqref{rw:anchors:signed-entry}, the whole
union has at most one. These statements concern the common physical
passing interval, with all whole-path and central positivity conditions.
\end{proposition}
\begin{proof}
The upper input $u$ decreases outward, so each positive-multiplier
numerator of $K_u$ has sign opposite to $K_p$. Throughout this proof
we restrict the quotients to the physical interval where the lower
whole-path test from $\chi(u)$, $u>\max(0,-\gamma)$, the finite
lower deficit and central positivity all hold.

\textit{The weak center.}
Write $d=d_s$ and $s=\varsigma=\bar B+\bar m$ in this proof.
Use the complete numerator \eqref{rw:infinity:weak-center-numerator},
with the actual moving eigenvalue $\gamma$ and the retained term
$-\epsilon F_I(\chi)A_U$ in $C$. On this fixed fibre,
Proposition~\ref{rw:coefficient-transport} supplies
$|\nu|\le\eta(|d|+|s|)$, $\eta\to0$.
With $e_s=e_\varsigma$, Lemma~\ref{rw:infinity:weak-entry-data}
supplies the divided coefficients
\begin{equation}\label{rw:anchors:weak-columns}
 M=d\overline C-t(se_s+\nu e_\nu),\qquad
 e_s>c,\quad \overline C'>1/2,
\end{equation}
on the same coefficient branch, with bounded $C_u^2$ coefficients
and small $\overline C$. Under that forcing the coordinate norm is
comparable to $|\tau_1|+|\sigma_1|$.

If $|s|<\sqrt\eta(|d|+|s|)$, the small-$s$ alternative of that
lemma gives $M/d$ strictly increasing, so $Z(M)\le1$.
Lemma~\ref{rw:anchors:calculus} gives all the assertions in this case.

Otherwise $s\ne0$ and
$e=e_s+(\nu/s)e_\nu>c/2$. If $d=0$, $M=-tse$ has no zero and
the assertions follow by the same Rolle inequalities. If $d\ne0$, set
\[
 R_0=\frac{\overline C}{te},\qquad
 M=dte(R_0-s/d).
\]
Lemma~\ref{rw:infinity:weak-entry-data} gives unimodality of $R_0$,
interval strict superlevel sets, and strictly increasing crossings
of nonpositive levels. In particular,
if $ds<0$, there is at most one zero of $M$, and the preceding
ordinary and away Rolle arguments apply.

It remains that $d$ and $s$ have the same sign. Since $\bar c>0$,
this is the sign of $\sigma_1$. Suppose first that the seam has
$K$ of sign $-\sigma_1$. Its regular derivative has that same sign,
so the seam numerator has sign $d$. Thus the seam lies in
\[
             \{u:R_0(u)>s/d\}.
\]
Moving outward decreases $u$ and can leave this interval at most once.
Consequently the remaining terminal part has at most one numerator
zero. Its away entrance gives at most one zero of $K$. This argument
applies after the possible regular zero and under the prescribed
signed entrance.

If instead the seam value of $K$ has sign $\sigma_1$, no regular
zero has occurred. If $M$ there has sign $d$, the same superlevel
argument gives at most one numerator zero and ordinary Rolle gives
at most two zeros of $K$. If $M$ has the opposite sign, then $K_p$
has the sign of $K$ at the seam; away Rolle and $Z(M)\le2$ again
give at most two. The seam avoids numerator zeros. If the original
overlap point was such a zero, its initial segment was assigned to $R$
before the seam was fixed, so the count still covers the entire union.

\textit{Positive $\bar B$.}
At a simple threshold $\bar a=-\bar c-\bar B/\bar c$ with
$0<\bar B<\bar c^2$, and at the double threshold
$\bar B=\bar c^2$, $\bar a=-2\bar c$. Hence at either threshold
$\tau_1<0$ and $\sigma_1>0$. Lemma~\ref{rw:infinity:weak-entry-data}
gives $K>0$ on the retained compact regular overlap for this same
original central function.

There was therefore no preceding regular zero, since its outward
sign would be $-\sigma_1<0$. For the same reason the prescribed
signed entry cannot occur on this union. It remains to bound the
unanchored terminal continuation of this positive overlap.
Use the numerator and quotient from
Proposition~\ref{rw:infinity:weak-loss},
\[
 M=AC-tE,\qquad E>c,\qquad R_0=C/(tE).
\]
Here $A=\delta_L$ for a simple lower loss, and
$A=\mu_L+\delta_L^2$ for a double loss.
Lemma~\ref{rw:infinity:weak-entry-data} gives $A>0$ and $C>0$
at every numerator zero and interval strict superlevel sets of $R_0$.
If $A\le0$, $M$ has no
zero and the assertion follows. If $A>0$, an entrance with $M>0$
lies in the interval $\{R_0>1/A\}$ and encounters at most one
numerator zero outward; ordinary Rolle gives at most two zeros of $K$.
An entrance with $M<0$ has $K_p>0$ and is away from zero, so
$Z(M)\le2$ and away Rolle give the same bound.

Reflection transports the union, its signs and multiplicities.
The improvement comes from this entrance; an arbitrary weak-chart
subinterval still has the unanchored four-zero closing estimate of
Section~8.
\end{proof}
